% resumo em português
\setlength{\absparsep}{18pt} % ajusta o espaçamento dos parágrafos do resumo

\begin{resumo}
    Esta monografia apresenta aplicação de análise de redes complexas a dados históricos olímpicos (1896-2021), modelando competições esportivas através de grafos direcionados e ponderados que revelam estruturas competitivas latentes. O trabalho analisa 8.679 atletas medalhistas em seis modalidades (Swimming, Basketball, Football, Athletics, Judo, Boxing) ao longo de 125 anos de história olímpica.

    As redes construídas representam atletas como nós e relações competitivas baseadas em confrontos diretos no pódio como arestas ponderadas por hierarquia de medalhas. Através de PageRank adaptado ao contexto esportivo, métricas de centralidade e detecção de comunidades via algoritmo de Louvain, o estudo identifica hierarquias de performance, agrupamentos estruturais e padrões temporais não evidentes em análises tradicionais baseadas em contagem de medalhas.

    Os resultados revelam diferenças sistemáticas entre modalidades: esportes individuais produzem redes altamente modulares com comunidades segregadas, enquanto esportes coletivos geram estruturas densamente interconectadas com baixa modularidade. A caracterização multidimensional de comunidades através de entropia temporal, concentração geográfica e índices de dominância permite compreensão profunda das dinâmicas competitivas olímpicas. Um dashboard interativo desenvolvido em Streamlit facilita exploração dos resultados e validação dos achados.

    \vspace{\onelineskip}

    \noindent
    \textbf{Palavras-chaves}: redes complexas, análise de desempenho esportivo, PageRank, detecção de comunidades, esporte olímpico, modelagem de grafos, métricas de centralidade, algoritmo de Louvain.
\end{resumo}
